\documentclass[graybox]{svmult}

\usepackage{mathptmx}
\usepackage{helvet}
\usepackage{courier}
\usepackage{type1cm}
\usepackage[T1]{fontenc}
\usepackage[utf8]{inputenc}
\usepackage{graphicx}
\usepackage{booktabs}
\usepackage{amsmath}
\usepackage{url}
\usepackage[hidelinks]{hyperref}
\usepackage{cite}
\usepackage[bottom]{footmisc}
\usepackage{microtype}

\newcommand{\nGauges}{19}
\newcommand{\nGaugesAlti}{7}
\newcommand{\nGaugesCB}{6}
\newcommand{\nGaugesWSA}{6}
\newcommand{\nPeru}{13}
\newcommand{\nChile}{6}
\newcommand{\nFeatures}{27}
\newcommand{\runtimeMin}{9}
\newcommand{\nBoot}{1000}
\newcommand{\bootBlock}{30}
\newcommand{\focusTestStart}{2021-09-22}
\newcommand{\focusTestEnd}{2024-12-24}
\newcommand{\focusTrainStart}{1994-09-25}
\newcommand{\focusNTestOne}{943}
\newcommand{\focusPersNseOne}{0.902}
\newcommand{\focusBestOne}{random forest}
\newcommand{\focusBestNseOne}{0.912}
\newcommand{\focusBestKgeOne}{0.924}
\newcommand{\focusBestSkillOne}{0.054}
\newcommand{\focusBestLoOne}{$-$0.003}
\newcommand{\focusBestHiOne}{0.113}
\newcommand{\focusBestDmpOne}{0.048}
\newcommand{\focusRidgeSkillOne}{0.070}
\newcommand{\focusNTestThree}{942}
\newcommand{\focusPersNseThree}{0.693}
\newcommand{\focusBestThree}{random forest}
\newcommand{\focusBestNseThree}{0.748}
\newcommand{\focusBestKgeThree}{0.822}
\newcommand{\focusBestSkillThree}{0.094}
\newcommand{\focusBestLoThree}{$-$0.004}
\newcommand{\focusBestHiThree}{0.160}
\newcommand{\focusBestDmpThree}{0.095}
\newcommand{\focusRidgeSkillThree}{0.100}
\newcommand{\focusNTestSeven}{941}
\newcommand{\focusPersNseSeven}{0.461}
\newcommand{\focusBestSeven}{Ridge regression}
\newcommand{\focusBestNseSeven}{0.578}
\newcommand{\focusBestKgeSeven}{0.677}
\newcommand{\focusBestSkillSeven}{0.115}
\newcommand{\focusBestLoSeven}{0.028}
\newcommand{\focusBestHiSeven}{0.208}
\newcommand{\focusBestDmpSeven}{0.112}
\newcommand{\focusRidgeSkillSeven}{0.115}
\newcommand{\nPosOne}{16}
\newcommand{\nSigOne}{11}
\newcommand{\nEvalOne}{19}
\newcommand{\medSkillOne}{0.051}
\newcommand{\medPersNseOne}{0.929}
\newcommand{\medBestNseOne}{0.934}
\newcommand{\medSkillAltiOne}{0.060}
\newcommand{\medSkillCBOne}{0.033}
\newcommand{\medSkillWSAOne}{0.053}
\newcommand{\nRidgePosOne}{14}
\newcommand{\nPosThree}{17}
\newcommand{\nSigThree}{12}
\newcommand{\nEvalThree}{19}
\newcommand{\medSkillThree}{0.080}
\newcommand{\medPersNseThree}{0.765}
\newcommand{\medBestNseThree}{0.787}
\newcommand{\medSkillAltiThree}{0.062}
\newcommand{\medSkillCBThree}{0.078}
\newcommand{\medSkillWSAThree}{0.083}
\newcommand{\nRidgePosThree}{18}
\newcommand{\nPosSeven}{17}
\newcommand{\nSigSeven}{12}
\newcommand{\nEvalSeven}{19}
\newcommand{\medSkillSeven}{0.108}
\newcommand{\medPersNseSeven}{0.546}
\newcommand{\medBestNseSeven}{0.620}
\newcommand{\medSkillAltiSeven}{0.103}
\newcommand{\medSkillCBSeven}{0.097}
\newcommand{\medSkillWSASeven}{0.130}
\newcommand{\nRidgePosSeven}{18}
\newcommand{\mostSelected}{Ridge regression}
\newcommand{\mostSelectedN}{11}
\newcommand{\pctDeltaTarget}{66}
\newcommand{\medCsiPers}{0.43}
\newcommand{\medCsiBest}{0.25}
\newcommand{\medCsiXgbLog}{0.33}
\newcommand{\medFarPers}{0.39}
\newcommand{\medFarBest}{0.38}
\newcommand{\medPodPers}{0.60}
\newcommand{\medPodBest}{0.44}
\newcommand{\nCsiBestBeats}{3}
\newcommand{\nCsiXgbLogBeats}{2}
\newcommand{\medCsiCalPers}{0.45}
\newcommand{\medCsiCalBest}{0.40}
\newcommand{\medPodCalBest}{0.68}
\newcommand{\medFarCalBest}{0.53}
\newcommand{\nCsiCalBestBeats}{3}
\newcommand{\nCsiCalImproves}{11}
\newcommand{\nFusion}{7}
\newcommand{\fusionMedDeltaRidge}{$-$0.005}
\newcommand{\fusionNPosRidge}{2}
\newcommand{\fusionMedDeltaXGBoost}{$-$0.008}
\newcommand{\fusionNPosXGBoost}{2}
\newcommand{\glofasMedNse}{$-$0.10}
\newcommand{\glofasNPosNse}{1}
\newcommand{\nChirps}{13}
\newcommand{\chirpsMedDeltaRidge}{$-$0.009}
\newcommand{\chirpsNPosRidge}{4}
\newcommand{\chirpsMaxDeltaRidge}{0.004}
\newcommand{\chirpsMedDeltaXGBoost}{$-$0.013}
\newcommand{\chirpsNPosXGBoost}{6}
\newcommand{\chirpsMaxDeltaXGBoost}{0.602}
\newcommand{\eraChirpsRatio}{2.8}
\newcommand{\wfNFolds}{57}
\newcommand{\wfNPos}{53}
\newcommand{\wfMedSkill}{0.103}
\newcommand{\wfFocusSkills}{0.166, 0.116, 0.114}
\newcommand{\nPairs}{104}
\newcommand{\corrDzero}{582}
\newcommand{\corrDzeroLo}{330}
\newcommand{\corrDzeroHi}{1069}
\newcommand{\corrC}{0.35}
\newcommand{\rhoDist}{$-$0.56}
\newcommand{\pDist}{$<$0.001}
\newcommand{\rhoSameRegion}{0.28}
\newcommand{\rhoCrossRegion}{0.09}
\newcommand{\moranI}{$-$0.012}
\newcommand{\moranP}{0.540}
\newcommand{\moranExp}{$-$0.056}
\newcommand{\topDriver}{snow fraction}
\newcommand{\topDriverRho}{$-$0.45}
\newcommand{\topDriverP}{0.056}
\newcommand{\topDriverPholm}{0.394}
\newcommand{\githubUrl}{https://github.com/Andre031222/streamflow-drought-altiplano}
\newcommand{\orcidAndre}{0009-0003-2385-5263}
\newcommand{\orcidMaribel}{0009-0003-6218-2735}
\newcommand{\orcidRenato}{0009-0005-2367-7852}
\newcommand{\emailAndre}{75521963@est.unap.edu.pe}
\newcommand{\emailMaribel}{maribelbel201314@gmail.com}
\newcommand{\emailRenato}{72535253@est.unap.edu.pe}

% Telephone of the corresponding author (required by Springer, not published).
\newcommand{\phoneAndre}{}
\InputIfFileExists{../private.tex}{}{}% private.tex (git-ignored) sets \phoneAndre for the CMT submission only
\newcommand{\phoneLine}{\ifx\phoneAndre\empty\else,\\ \hphantom{\textbf{Corresponding author:}} Tel.:~\mbox{\phoneAndre}\fi}
\makeatletter
\newcommand{\inputtable}[1]{\@@input #1 }% primitive input: safe inside tabular
\makeatother

\begin{document}

\title*{Honest Baselines for Streamflow and Hydrological-Drought Forecasting in the Tropical Andes: An Open, Leakage-Free Multi-Station Benchmark with an Altiplano Focus}
\titlerunning{Honest Andean Streamflow and Drought Forecasting Benchmark}
\author{Richar Andre Vilca-Solorzano, Dina Maribel Yana-Yucra and Renato Quispe-Vargas}
\authorrunning{R. A. Vilca-Solorzano et al.}
\maketitle

\vspace*{-26mm}
\noindent{\small Faculty of Statistical and Computer Engineering, Universidad Nacional del Altiplano, Av.~Floral 1153, Puno 21001, Peru\\[3pt]
\textbf{Corresponding author:} Richar Andre Vilca-Solorzano, e-mail: \href{mailto:\emailAndre}{\emailAndre}\phoneLine\\
Dina Maribel Yana-Yucra, e-mail: \href{mailto:\emailMaribel}{\emailMaribel}\\
Renato Quispe-Vargas, e-mail: \href{mailto:\emailRenato}{\emailRenato}\\[3pt]
ORCID: \orcidAndre{} (R.\,A.\,V.-S.), \orcidMaribel{} (D.\,M.\,Y.-Y.), \orcidRenato{} (R.\,Q.-V.)}
\par\bigskip

\abstract*{Short-range streamflow forecasts and early drought signals matter in the tropical Andes, yet reported machine-learning gains are hard to compare because evaluation is rarely free of information leakage. We release an open benchmark of \nGauges{} long-record gauges from the Andean Glacierized Catchment dataset: \nGaugesAlti{} on the Altiplano, including Puente Isla Cabanillas in the Lake Titicaca basin, \nGaugesCB{} in the Cordillera Blanca and \nGaugesWSA{} in the western and southern Andes. The gauges are combined with ERA5 and MERRA-2 reanalysis forcing and GloFAS~v4 simulated discharge. Persistence and climatology are compared with five families of learners at 1-, 3- and 7-day leads, under a calendar-continuous, embargoed and validation-tuned protocol with block-bootstrap confidence intervals, Diebold--Mariano tests and walk-forward refits. The learners beat persistence at \nPosOne{} of \nGauges{} gauges at one day and at \nPosSeven{} of \nGauges{} at seven days, with a median skill of \medSkillSeven{}. At Cabanillas, the one-day skill is \focusBestSkillOne{} (95\% confidence interval \focusBestLoOne--\focusBestHiOne), which reverses an earlier result obtained with a leaky pipeline. A lower root-mean-square error did not improve the detection of drought days: the median critical success index is \medCsiPers{} for persistence and \medCsiBest{} for the learners. Hydrological drought decorrelates over about \corrDzero{}~km, and forecast skill follows the hydroclimatic regime rather than location.}

\abstract{Short-range streamflow forecasts and early drought signals matter in the tropical Andes, yet reported machine-learning gains are hard to compare because evaluation is rarely free of information leakage. We release an open benchmark of \nGauges{} long-record gauges from the Andean Glacierized Catchment dataset: \nGaugesAlti{} on the Altiplano, including Puente Isla Cabanillas in the Lake Titicaca basin, \nGaugesCB{} in the Cordillera Blanca and \nGaugesWSA{} in the western and southern Andes. The gauges are combined with ERA5 and MERRA-2 reanalysis forcing and GloFAS~v4 simulated discharge. Persistence and climatology are compared with five families of learners at 1-, 3- and 7-day leads, under a calendar-continuous, embargoed and validation-tuned protocol with block-bootstrap confidence intervals, Diebold--Mariano tests and walk-forward refits. The learners beat persistence at \nPosOne{} of \nGauges{} gauges at one day and at \nPosSeven{} of \nGauges{} at seven days, with a median skill of \medSkillSeven{}. At Cabanillas, the one-day skill is \focusBestSkillOne{} (95\% confidence interval \focusBestLoOne--\focusBestHiOne), which reverses an earlier result obtained with a leaky pipeline. A lower root-mean-square error did not improve the detection of drought days: the median critical success index is \medCsiPers{} for persistence and \medCsiBest{} for the learners. Hydrological drought decorrelates over about \corrDzero{}~km, and forecast skill follows the hydroclimatic regime rather than location.
\keywords{Streamflow forecasting $\cdot$ Hydrological drought $\cdot$ Altiplano $\cdot$ Machine learning $\cdot$ Persistence baseline $\cdot$ Data leakage $\cdot$ Spatial statistics $\cdot$ Open benchmark}}

\section{Introduction}
\label{sec:intro}

Streamflow in the tropical Andes sustains irrigation, hydropower and drinking water for millions of people, yet it is shaped by an extreme seasonal cycle, retreating glaciers and recurrent drought \cite{vuille2018andes,baraer2012glacier,lavado2012peru}. On the Altiplano, the endorheic plateau shared by Peru, Bolivia and Chile, rivers such as the Coata--Cabanillas feed Lake Titicaca and depend almost entirely on a short austral-summer wet season \cite{garreaud2003altiplano}; climate projections point to more frequent and longer hydrological droughts in the Titicaca--Desaguadero--Poop\'o system \cite{zubieta2021titicaca}. Short-range discharge forecasts and early drought signals are therefore valuable. In practice, however, gauges are sparse, records are gappy and operational tools are rarely open.

Machine learning (ML) is increasingly used for data-driven streamflow forecasting \cite{yaseen2015ai,mosavi2018flood,kratzert2018lstm,kratzert2019lstm,nearing2021role}. For daily lead times, however, discharge is strongly autocorrelated, and the naive persistence forecast $\hat{Q}_{t+h}=Q_t$ is a demanding benchmark \cite{pappenberger2015benchmarks,murphy1988skill}. Reported gains can be inflated by (i) comparing models against weak references, (ii) randomly shuffled cross-validation that leaks temporal information \cite{bergmeir2012cv,roberts2017cv,torres2026blockedcv}, and (iii) implementation details that give baselines and learners different information at issue time. An earlier version of our own pipeline suffered from the third problem: the learners saw discharge only up to $t-1$ while persistence used $Q_t$, and calendar gaps were treated as consecutive days. Correcting both reversed the one-day conclusion, and this experience motivated the explicit protocol reported here.

This paper contributes an open, laptop-scale benchmark for the Peruvian Andes and the Altiplano in which every evaluation choice is explicit and tested. Specifically, we:
\begin{enumerate}
\item assemble \nGauges{} long-record gauges from the Andean-GC dataset \cite{aguayo2026andeangc}, covering three regions: \nGaugesAlti{} Altiplano gauges (Cabanillas, Peru, and the Lauca/Caquena headwaters, Chile), \nGaugesCB{} glacier-fed Cordillera Blanca gauges and \nGaugesWSA{} western/southern Andean gauges. We combine them with forcing from ERA5, the fifth-generation reanalysis of the European Centre for Medium-Range Weather Forecasts (ECMWF) \cite{hersbach2020era5}; with precipitation from NASA's Prediction of Worldwide Energy Resources (POWER) service, which is based on the Modern-Era Retrospective analysis for Research and Applications, version~2 (MERRA-2); and with discharge simulated by the Global Flood Awareness System (GloFAS)~v4;
\item define a leakage-free protocol: a calendar-continuous index, identical information at issue time for all methods, chronological splits with an embargo, hyper-parameter selection on validation only, and expanding-window walk-forward checks;
\item report skill against persistence with moving-block bootstrap intervals \cite{kunsch1989block} and Diebold--Mariano tests \cite{diebold1995dm,harvey1997dm}, at 1-, 3- and 7-day leads;
\item evaluate drought detection directly, using the variable-threshold level method \cite{vanloon2015explained,fleig2006threshold} and categorical scores \cite{schaefer1990csi}, and release a nonparametric standardized streamflow index (SSI) \cite{farahmand2015nonparametric}.
\end{enumerate}
The paper is organized as follows. Section~\ref{sec:related} reviews related work, Sect.~\ref{sec:data} describes the data and Sect.~\ref{sec:methods} the methods. Section~\ref{sec:results} presents the results, Sect.~\ref{sec:discussion} discusses them and Sect.~\ref{sec:conclusion} concludes.

\section{Related Work}
\label{sec:related}

\textbf{Data-driven streamflow forecasting in the Andes.} Random forests have been used for runoff and flash-flood forecasting in Ecuadorian Andean catchments \cite{munoz2018rf}, recently combined with satellite precipitation and hydrological knowledge \cite{munoz2025rf}. Support-vector regression was tested for monthly runoff in the Central Andes of Argentina \cite{teverovsky2023svr}. In Peru, artificial neural networks have forecast monthly flows in the high-Andean Crisnejas basin \cite{ramos2023crisnejas}, and $k$-nearest neighbors daily flows of the Ramis River on the Altiplano \cite{quispe2023ramis}. Ensemble data assimilation with a conceptual model improved hourly forecasts in the Vilcanota basin \cite{llauca2023vilcanota}. On the Altiplano, Lujano et al.\ found that reanalysis precipitation performs well against gauges, but gridded forcing can still degrade conceptual-model streamflow \cite{lujano2025gridded}. This result motivates both our reanalysis-based features and the data-fusion ablation. Large-sample studies with long short-term memory (LSTM) networks set the methodological reference for data-driven hydrology \cite{kratzert2018lstm,kratzert2019lstm,gauch2021lstm,addor2017camels,nearing2021role}. Most Andean studies, however, evaluate a single basin at a single lead time, using NSE without a persistence reference or an uncertainty estimate. Our prior multi-station climate-ML study of Altiplano frost \cite{torres2025frost} used the same comparative-benchmark design.

\textbf{Validation under dependence.} Random $k$-fold cross-validation is optimistic for autocorrelated or spatially structured data \cite{bergmeir2012cv,roberts2017cv}. For Peruvian agricultural survey data, blocked cross-validation revealed that random folds overstated accuracy \cite{torres2026blockedcv}, and hierarchical spatial models were needed to characterize climate vulnerability \cite{torres2026bayescar}. We transfer the same principle to time: chronological blocks with an embargo, and walk-forward refits.

\textbf{Hydrological drought.} Hydrological drought is a below-normal availability of streamflow or storage \cite{vanloon2015explained,mishra2010review}. It is identified with threshold-level methods \cite{fleig2006threshold,vanloon2015severity} or with standardized runoff and streamflow indices (SRI, SSI) \cite{shukla2008sri,vicenteserrano2012ssi}. These are analogues of the Standardized Precipitation Index (SPI) and the Standardized Precipitation Evapotranspiration Index (SPEI) \cite{mckee1993spi,vicenteserrano2010spei}. Parametric SSI values are sensitive to the chosen distribution \cite{tijdeman2020ssi}, which motivates nonparametric plotting-position formulations \cite{farahmand2015nonparametric,gringorten1963}. Drought forecasting requires an explicit identification step after the streamflow forecast, and the chosen definition changes the result \cite{sutanto2021definitions,hao2018prediction,sucozhanay2026threshold}. We therefore score drought days explicitly instead of inferring them from RMSE.

\begin{figure}[!t]
\centering
\includegraphics[width=0.88\textwidth]{figuras/fig00_station_map.pdf}
\caption{The \nGauges{} benchmark gauges and their Andean-GC catchments, with marker shape and shading indicating the region. The Altiplano focus gauge Puente Isla Cabanillas is highlighted. Insets show the Cordillera Blanca and Chilean Altiplano clusters.}
\label{fig:map}
\end{figure}

\section{Study Area and Data}
\label{sec:data}

\textbf{Discharge.} We use quality-controlled daily discharge from the Andean-GC dataset \cite{aguayo2026andeangc}, released under CC~BY~4.0. We keep every Peruvian gauge, operated by the National Meteorological and Hydrological Service (SENAMHI), and every Altiplano gauge of the Lauca/Caquena headwaters in Chile, operated by the General Directorate of Water (DGA), with at least 5000 quality-controlled (QC) days, which gives \nGauges{} gauges (\nPeru{} Peru, \nChile{} Chile; Table~\ref{tab:gauges}, Fig.~\ref{fig:map}). The focus gauge is Puente Isla Cabanillas (\texttt{P00000055}; 2843~km$^2$, mean elevation 4471~m), which drains the Coata--Cabanillas system into Lake Titicaca. The Cordillera Blanca gauges are historical records from 1952--1974 in strongly glacierized catchments. The remaining gauges span 1950--2024. Records contain multi-year gaps (e.g., 1996--2005 at Cabanillas), which we keep as missing values on a daily calendar.

\textbf{Meteorological forcing.} Daily precipitation and 2-m temperature come from ERA5 \cite{hersbach2020era5} at each catchment centroid, retrieved through the Open-Meteo historical API \cite{openmeteo}. The API's ERA5-Land endpoint \cite{munozsabater2021era5land} does not provide precipitation, so we use ERA5 throughout. A second, independent precipitation estimate is NASA POWER \texttt{PRECTOTCORR} (MERRA-2 bias-corrected with GPCP; available from 1981).

\textbf{Simulated discharge.} GloFAS~v4 reanalysis discharge is retrieved from the Open-Meteo Flood API (available from 1984) on a $3\times3$ neighborhood of 0.05$^\circ$ cells around each gauge. The cell used is the one with the highest Kling--Gupta efficiency (KGE) over the fitting period only, so the test period is never used for this choice.



\begin{table}[!htb]
\centering
\caption{Benchmark gauges (Andean-GC, $\ge$5000 QC days). Area in km$^2$; $n_{\mathrm{test}}$ is the number of test days at a 7-day lead. $^\star$Focus gauge.}
\label{tab:gauges}
\scriptsize
\setlength{\tabcolsep}{3pt}
\begin{tabular}{llllrlrr}
\toprule
ID & Gauge & Cty & Region & Area & Record & QC days & $n_{\mathrm{test}}$ \\
\midrule
\inputtable{tab_gauges.tex}
\bottomrule
\end{tabular}
\end{table}

\section{Methods}
\label{sec:methods}

\textbf{Forecast set-up.} For each gauge and lead $h\in\{1,3,7\}$ days, a forecast issued at day $t$ predicts $Q_{t+h}$ using information available at $t$. The \nFeatures{} predictors are: discharge lags $Q_{t-k}$, $k\in\{0,1,2,3,7,14,30\}$; the one-day change; rolling means and standard deviations over 7, 14 and 30 days ending at $t$; precipitation lags (0--7 days) and 7-, 30- and 90-day totals; 2-m temperature (daily, 7- and 30-day means); and the day of year encoded as sine and cosine. All operations use a daily calendar index, so a lag of $k$ days is exactly $k$ days. Any sample whose predictors or target fall in a data gap is dropped. Persistence uses the same $Q_t$ that the learners see.

\textbf{Methods compared.} There are two references: persistence and a day-of-year climatology (15-day smoothed mean of the fitting period, evaluated at the target date). We compare six learners: Ridge regression \cite{hoerl1970ridge}, random forest \cite{breiman2001rf}, histogram gradient boosting \cite{friedman2001gbm}, XGBoost \cite{chen2016xgboost}, a two-layer multilayer perceptron (MLP), and a low-flow-oriented XGBoost trained on the change in $\log(1+Q)$ (XGBoost-log). For every learner, a small grid of hyper-parameters is searched together with the target form: the level $Q_{t+h}$ or the change $Q_{t+h}-Q_t$.

\textbf{Validation protocol.} Each record is split chronologically into 70\% training, 15\% validation and 15\% test. The last $h$ samples of the training and validation blocks are removed (embargo) so that no target overlaps the next block. Configurations are selected on validation RMSE. The selected configuration is refit on training plus validation and scored once on the test block. When we refer to the learner of a gauge, we mean the one with the lowest validation RMSE, so the test set is never used to choose models. As a robustness check, an expanding-window walk-forward with three folds is run for Ridge at every gauge, re-selecting hyper-parameters within each fold.

\textbf{Scores and uncertainty.} We report the root-mean-square error (RMSE), the Nash--Sutcliffe efficiency (NSE) \cite{nash1970}, the KGE \cite{klinggupta2009,knoben2019kge} and the RMSE skill against persistence, $S=1-\mathrm{RMSE}_{\mathrm{model}}/\mathrm{RMSE}_{\mathrm{pers}}$ \cite{murphy1988skill}. 95\% confidence intervals (CIs) for $S$ come from a paired moving-block bootstrap (\nBoot{} replicates, \bootBlock-day blocks) that preserves autocorrelation \cite{kunsch1989block,politis1994stationary}. Equal squared-error accuracy against persistence is tested with the Diebold--Mariano test using the small-sample correction \cite{diebold1995dm,harvey1997dm}.

\textbf{Drought detection.} Following the variable-threshold level method \cite{vanloon2015explained,fleig2006threshold,sutanto2021definitions}, a day is in hydrological drought when $Q$ falls below the monthly 20th percentile (Q80 exceedance) of the fitting period. For each method, we compare the forecast drought state at $t+h$ with the observed state and compute the probability of detection (POD), false-alarm ratio (FAR) and critical success index (CSI) \cite{schaefer1990csi}. For description we also compute a nonparametric SSI-30: the 30-day mean discharge is ranked within each calendar month, converted to a Gringorten plotting position, and mapped to a standard-normal quantile \cite{farahmand2015nonparametric,gringorten1963}. This avoids the distributional sensitivity of parametric SSI \cite{tijdeman2020ssi}.

\textbf{Data-fusion ablation.} Where MERRA-2 precipitation and GloFAS overlap the record (1984 onwards), we refit Ridge and XGBoost on identical samples with two predictor sets: ERA5 only, and ERA5 plus MERRA-2 precipitation and lagged GloFAS discharge up to $t$. GloFAS values at the target date are never used as predictors. We also report the raw GloFAS simulation as a physically based reference.

\textbf{Spatial statistics.} Spatial drought coherence is quantified as the Spearman correlation of monthly SSI-30 between every gauge pair with at least 60 overlapping months, as a function of great-circle distance $d$. We fit the exponential correlogram $\rho(d)=c\,e^{-d/d_0}$, in the spirit of space--time drought analysis \cite{tallaksen2009spacetime}. Pairs share gauges and are therefore not independent, so the 95\% interval of the decorrelation distance $d_0$ comes from a gauge-level bootstrap (500 replicates). Spatial clustering of the 7-day skill is tested with global Moran's $I$ \cite{moran1950}, using inverse-distance weights and 9999 permutations. Candidate drivers of skill are Andean-GC catchment attributes (area, mean elevation, glacier fraction, aridity, precipitation concentration, snow fraction and dry-day frequency). They are screened with Spearman correlations, permutation $p$-values and a Holm correction for multiple testing. The skill map and the drought-synchrony network are rendered in QGIS.

The whole pipeline runs on a laptop CPU (10 cores, about \runtimeMin{}~min), and every number in this paper is written automatically from the result tables, so text and results cannot drift apart.

\section{Results}
\label{sec:results}

\subsection{Focus gauge: Puente Isla Cabanillas}
Figure~\ref{fig:focus} shows the Cabanillas record, the chronological blocks and the nonparametric SSI-30. Table~\ref{tab:focus} reports all methods on the test period (\focusTestStart{} to \focusTestEnd{}). With identical information at issue time, persistence is no longer unbeatable at one day. The validation-selected \focusBestOne{} reaches NSE~$=$~\focusBestNseOne{} against \focusPersNseOne{} for persistence, a skill of \focusBestSkillOne{} (95\% CI \focusBestLoOne{} to \focusBestHiOne{}; Diebold--Mariano $p$~\focusBestDmpOne{}). This reverses the one-day conclusion of the earlier, leaky version of the pipeline (Sect.~\ref{sec:intro}). At seven days, persistence drops to NSE~$=$~\focusPersNseSeven{}. Ridge regression on the change target gives the most robust gain (skill \focusRidgeSkillSeven{}, CI excluding zero), and its walk-forward folds give skills of \wfFocusSkills{}. The learner selected on validation at this lead (\focusBestSeven{}) has a smaller, non-significant skill of \focusBestSkillSeven{}. Validation-based selection is unbiased, but it is noisy when validation and test hydrology differ.

\begin{figure}[t]
\centering
\includegraphics[width=\textwidth]{figuras/fig01_focus_hydrograph_ssi.pdf}
\caption{Puente Isla Cabanillas: daily discharge (log scale) with the chronological train/validation/test blocks, and the nonparametric SSI-30 (bars below zero: below-normal flow). Gaps are real missing records.}
\label{fig:focus}
\end{figure}

\begin{table}[t]
\centering
\caption{Test-period scores at Puente Isla Cabanillas. Skill is relative to persistence with 95\% block-bootstrap CI. DM: Diebold--Mariano $p$-value. CSI: drought-day detection. $^\dagger$Learner selected on validation RMSE.}
\label{tab:focus}
\scriptsize
\setlength{\tabcolsep}{3.2pt}
\begin{tabular}{rlrrrrcrr}
\toprule
$h$ & Method & RMSE & NSE & KGE & Skill & 95\% CI & DM $p$ & CSI \\
\midrule
\inputtable{tab_focus.tex}
\bottomrule
\end{tabular}
\end{table}

\subsection{Multi-station benchmark}
Across the \nGauges{} gauges, the validation-selected learner beats persistence at \nPosOne{}/\nEvalOne{} gauges at one day, \nPosThree{}/\nEvalThree{} at three days and \nPosSeven{}/\nEvalSeven{} at seven days. Its 95\% interval excludes zero at \nSigOne{}, \nSigThree{} and \nSigSeven{} gauges, respectively (Fig.~\ref{fig:heatmap}). The median skill grows with lead time, from \medSkillOne{} to \medSkillThree{} and \medSkillSeven{}, because persistence degrades faster than the learners (Fig.~\ref{fig:lead}). The regions behave differently. Skill rises steeply with lead time in the Cordillera Blanca (median \medSkillCBOne{} to \medSkillCBSeven{}) and in the western/southern Andes (\medSkillWSAOne{} to \medSkillWSASeven{}). On the Altiplano it stays flat (\medSkillAltiOne{} to \medSkillAltiSeven{}), and the smallest Chilean headwater (Caquena at its source, 50~km$^2$) remains negative at seven days. Model choice matters less than the protocol: Ridge alone is skilful at \nRidgePosSeven{}/\nGauges{} gauges at seven days. The change target was selected in \pctDeltaTarget\% of the learner fits. The expanding-window walk-forward confirms the pattern: \wfNPos{} of \wfNFolds{} gauge--fold combinations have positive seven-day Ridge skill (median \wfMedSkill{}).

\begin{figure}[t]
\centering
\includegraphics[width=0.78\textwidth]{figuras/fig03_skill_heatmap_h7.pdf}
\caption{Seven-day RMSE skill against persistence for every gauge and method. Boxed cells: 95\% moving-block bootstrap interval above zero. Gauges are grouped by region (top to bottom: Altiplano, Cordillera Blanca, western/southern Andes).}
\label{fig:heatmap}
\end{figure}

\begin{figure}[t]
\centering
\includegraphics[width=0.9\textwidth]{figuras/fig04_lead_skill_by_region.pdf}
\caption{Skill of the validation-selected learner versus lead time, by region (thin lines: gauges; thick: median).}
\label{fig:lead}
\end{figure}

\subsection{Hydrological-drought detection}
Lower RMSE does not translate into better drought detection (Fig.~\ref{fig:drought}). At seven days, the median CSI of drought days is \medCsiPers{} for persistence and \medCsiBest{} for the validation-selected learner. The learner beats persistence at only \nCsiBestBeats{}/\nGauges{} gauges. The low-flow-oriented XGBoost-log does not close the gap (median CSI \medCsiXgbLog{}; better than persistence at \nCsiXgbLogBeats{} gauges). The learners are more conservative: they raise fewer false alarms (median FAR \medFarBest{} vs.\ \medFarPers{}) but miss more drought days (POD \medPodBest{} vs.\ \medPodPers{}). The reason is that RMSE-trained models regress towards the conditional mean and smooth the recession limb, near where the threshold lies. This agrees with the view that drought forecasting needs its own identification and verification step \cite{sutanto2021definitions}.

\begin{figure}[t]
\centering
\includegraphics[width=\textwidth]{figuras/fig05_drought_detection_h7.pdf}
\caption{Seven-day detection of drought days (variable threshold: monthly Q20 of the fitting period). Left: CSI; right: false-alarm ratio.}
\label{fig:drought}
\end{figure}

\subsection{Data fusion and the GloFAS reference}
At the \nFusion{} gauges where MERRA-2 and GloFAS overlap the record, adding them to the ERA5 predictors changes the seven-day skill by a median of \fusionMedDeltaRidge{} for Ridge and \fusionMedDeltaXGBoost{} for XGBoost. Only \fusionNPosRidge{} and \fusionNPosXGBoost{} gauges improve, respectively. The raw GloFAS~v4 simulation, even with the cell chosen on the fitting period, has a median NSE of \glofasMedNse{} (positive at \glofasNPosNse{}/\nFusion{} gauges). At these small, steep Andean catchments, recent observed discharge carries far more short-range information than a global model or a second reanalysis.

\subsection{Spatial statistics}
Hydrological drought is spatially coherent but local (Figs.~\ref{fig:spatialmap}b and~\ref{fig:spatial}a). Across \nPairs{} gauge pairs, the monthly SSI-30 correlation decreases with distance (Spearman \rhoDist{}, $p$~\pDist{}). It follows $\rho(d)=\corrC\,e^{-d/d_0}$ with a decorrelation distance $d_0=\corrDzero$~km (95\% CI \corrDzeroLo--\corrDzeroHi~km). Pairs within the same region correlate more strongly than pairs across regions (median $\rho$ \rhoSameRegion{} vs.\ \rhoCrossRegion{}). Even nearby gauges share only part of their drought signal, which shows that one gauge cannot represent drought on the Altiplano or in the Cordillera Blanca. Forecast skill, in contrast, shows no spatial clustering: Moran's $I=\moranI$ ($p=\moranP$; expectation under randomness \moranExp; Fig.~\ref{fig:spatialmap}a). Skill is better explained by hydroclimatic regime than by location (Fig.~\ref{fig:spatial}b). The strongest correlate is \topDriver{} ($\rho=\topDriverRho$, permutation $p=\topDriverP$, Holm-adjusted $p=\topDriverPholm$), followed by precipitation concentration and aridity. Skill tends to be lower in catchments with more snow, more concentrated precipitation and higher aridity, which is the hydroclimate of the Altiplano headwaters. However, only snow fraction approaches significance after correction for multiple testing.

\begin{figure}[t]
\centering
\includegraphics[width=\textwidth]{figuras/fig08_spatial_skill_synchrony.pdf}
\caption{QGIS maps of (a) seven-day skill of the validation-selected learner by catchment (bold outline: 95\% CI above zero) and (b) the hydrological-drought synchrony network (monthly SSI-30 Spearman correlation between gauge pairs with $\ge$60 overlapping months).}
\label{fig:spatialmap}
\end{figure}

\begin{figure}[t]
\centering
\includegraphics[width=0.92\textwidth]{figuras/fig09_spatial_statistics.pdf}
\caption{(a) Drought-synchrony correlogram with exponential fit (dotted: 95\% gauge-bootstrap range of $d_0$). (b) Spearman correlation between seven-day skill and catchment attributes (highlighted bars: permutation $p<0.05$, uncorrected).}
\label{fig:spatial}
\end{figure}

\section{Discussion}
\label{sec:discussion}

\textbf{Protocol before models.} The largest change in our conclusions came from the protocol, not from model choice. Giving persistence and learners the same information at issue time, and indexing on the calendar, turned a reported one-day loss into a significant gain. Hyper-parameter selection on validation and bootstrap intervals then separated robust gains from noise. We recommend that Andean ML forecasting studies report skill against persistence with uncertainty, publish their split and embargo design, and test leakage explicitly; our test suite does this automatically.

\textbf{Where ML helps.} Gains are modest, a median RMSE reduction of about 5--10\%, but consistent: positive at almost every gauge at seven days and confirmed by walk-forward folds. Linear models on the change target are competitive with tree ensembles, in line with large-sample findings that careful baselines are hard to beat at short leads \cite{pappenberger2015benchmarks,nearing2021role}. The Altiplano is the hardest region. Its headwaters combine short convective wet-season pulses with long baseflow recessions \cite{garreaud2003altiplano}. A 0.25$^\circ$ reanalysis can represent this rainfall only partially, and even well-performing gridded products degrade conceptual streamflow simulations there \cite{lujano2025gridded}.

\textbf{Drought is a different target.} Better RMSE did not yield better drought-day detection. Operational drought early warning should train or calibrate on the drought decision itself, for example with quantile or classification losses and threshold-aware post-processing, and verify it with categorical scores \cite{sutanto2021definitions,hao2018prediction}. The limited spatial coherence of drought ($d_0$ of a few hundred kilometres) argues for dense gauge networks rather than extrapolating from single stations, consistent with projected changes in the Titicaca system \cite{zubieta2021titicaca}.

\textbf{Limitations.} ERA5 at 0.25$^\circ$ is coarse for steep catchments and has a wet bias over the Andes. The Cordillera Blanca records predate 1975, so no MERRA-2 or GloFAS data overlap them. The fusion ablation therefore covers \nFusion{} gauges only, and some GloFAS requests were rate-limited. The validation block is a single period, so selection can be unstable under non-stationarity (e.g., glacier retreat). Forecasts use observed rather than forecast forcing beyond $t$, which is acceptable here because predictors stop at $t$; skill against operational meteorological forecasts is left for future work.

\section{Conclusion}
\label{sec:conclusion}

We released an open, leakage-free benchmark for short-range streamflow and hydrological drought forecasting at \nGauges{} Andean gauges with an Altiplano focus. Under an honest protocol, data-driven models beat persistence at most gauges and leads: at seven days at \nPosSeven{}/\nGauges{} gauges (median skill \medSkillSeven{}), and also at one day at the Altiplano focus gauge, Cabanillas. RMSE gains do not carry over to drought-day detection, where persistence remains a strong reference. Hydrological drought decorrelates over about \corrDzero~km, and forecast skill depends on hydroclimatic regime rather than location. Future work will add drought-aware training losses, operational meteorological forecasts as inputs, and more Altiplano gauges from SENAMHI and the National Water Authority (ANA).


\section*{Data and code availability}
The discharge data are from Andean-GC (doi:\href{https://doi.org/10.5281/zenodo.18035801}{10.5281/zenodo.18035801}; CC~BY~4.0). ERA5 and GloFAS were retrieved through Open-Meteo \cite{openmeteo}, and MERRA-2 precipitation through NASA POWER. Source code, data extracts, leakage tests, QGIS map scripts, result tables and supplementary figures are public at:\\ {\small\url{\githubUrl}}

\section*{Declarations}
\noindent\textbf{Author contributions.} R.\,A. Vilca-Solorzano conceived the study, developed the software and the analysis, and wrote the manuscript. D.\,M. Yana-Yucra contributed to the methodology, the analysis and the writing. R. Quispe-Vargas validated the results. All authors approved the final manuscript.

\smallskip\noindent\textbf{Competing interests.} The authors declare no competing interests.

\smallskip\noindent\textbf{Ethics and intended use.} No human subjects or personal data were involved. The forecasts are research benchmarks and do not replace the official products of SENAMHI, DGA or other national services.

\smallskip\noindent\textbf{Acknowledgements.} We thank the Andean-GC team and the hydrological services whose records it compiles (SENAMHI, Peru; DGA, Chile), Copernicus/ECMWF, NASA POWER and Open-Meteo for open data, and the Faculty of Statistical and Computer Engineering, Universidad Nacional del Altiplano, for academic support.


\bibliographystyle{spmpsci}
\bibliography{references}

\end{document}
